\documentclass[11pt]{article}
\usepackage[margin=0.75in]{geometry}
\usepackage{fontspec}
\usepackage{setspace}
\usepackage{hyperref}
\hypersetup{colorlinks=true, urlcolor=blue, linkcolor=black}
\pagestyle{empty}
\setlength{\parskip}{0.62em}
\setlength{\parindent}{0pt}
\setstretch{0.92}
\setmainfont{Times New Roman}

\begin{document}

\begin{flushleft}
Ke Cheng HSU (Clio) \\
Department of Economics, The Hong Kong University of Science and Technology \\
Clear Water Bay, Kowloon, Hong Kong $\cdot$ +852 6840 9753 \\
\href{mailto:khsu@connect.ust.hk}{khsu@connect.ust.hk} $\cdot$ \href{https://cliohsu1997.github.io/}{cliohsu1997.github.io}
\end{flushleft}

24 November 2026

\vspace{0.08em}
Hiring Committee \\
Lee Kong Chian School of Business \\
Singapore Management University \\
Singapore

\vspace{0.08em}
Dear Members of the Hiring Committee,

I am applying for the Research Fellow (1-year contract) position at LKCSB. I am finishing my Ph.D. in Economics at HKUST (expected June 2027) and can start on or after 4 January 2027. I work on empirical industrial organization and digital economics, which means I spend most of my time on large data pipelines and empirical models---both structural and reduced form.

Much of what you list under duties is what I already do. For my job market paper, I assembled Taobao apparel transactions from scraped product pages and OCR-extracted product images, then estimated a structural model of purchase, return-shipping-insurance take-up, and returns across six categories. The estimates show that automatic coverage (RI) works better in high-margin, high-salvage categories because it expands valuable purchases, while optional coverage (CI) works better in thin-margin categories by curbing excessive returns; a partial restocking-fee menu often dominates both. In ongoing work with Kohei Kawaguchi and Kei Kawai, I link high-frequency bookstore point-of-sale records (deliveries, sales, returns) with scraped title attributes and local demographics. I write the code in Python and R (with C++ where the simulation is heavy), and I use SQL-style joins to build these panels---reproducibly, on AWS and the university HPC cluster.

On methods, I work with the reduced-form toolkit---difference-in-differences, RDD, IV, field experiments---as well as structural estimation. My own papers implement GMM and IV estimators with bootstrap inference and Monte Carlo checks, then use counterfactual simulation to rank policies. I read literature quickly and turn it into papers, slides, and short reports on deadline.

You will also want someone who communicates and collaborates well. I work closely with Pak Hung Au (on product design and search costs) and with Kawaguchi and Kawai, and I have presented at APIOC, ESAM, AMES, and the China Economics Annual Conference, as well as HKUST workshops. I have also served as Teaching Assistant Coordinator and taught MSc Math Camp (4.53/5), with an average TA rating of 4.72/5.

I have uploaded my CV and this letter through the SMU careers portal.\\
Happy to discuss the project: \href{mailto:khsu@connect.ust.hk}{khsu@connect.ust.hk} or +852~6840~9753.

\vspace{0.18em}
Sincerely,

\vspace{0.18em}
Ke Cheng HSU \\
Ph.D.\ Candidate, Department of Economics, HKUST

\end{document}
